\section{Discussion}
\label{sec:discussion}

The results indicate selective rather than uniform changes across training regimes. DD-LoRA already alters ACL Bench performance substantially, but the direction and magnitude depend on the underlying checkpoint. RC-LoRA is higher than DD-LoRA in overall ACL Bench accuracy for all three evaluated checkpoints, while the category-level effects remain heterogeneous. For Qwen2.5-7B, RC-LoRA also exhibits a larger BCDM than Base and a different pattern of causal dependence across a limited subset of attention heads. These effects are related, but they should not be collapsed into a single claim of general capability improvement.

\subsection{Selective and Model-Dependent Transfer}
\label{sec:discussion_transfer}

The three-way ACL Bench comparison separates two effects that were conflated in a Base-versus-RC-LoRA analysis. First, ordinary structured-task adaptation already produces substantial transfer in some settings. DeepSeek-R1-Distill-Llama-8B increases from 28.3\% to 31.7\% overall under DD-LoRA, and Qwen2.5-7B increases from 35.3\% to 54.0\%, whereas Llama3.1-8B decreases from 45.0\% to 42.3\%. Ordinary Dou Dizhu adaptation therefore does not provide a uniformly beneficial transfer signal. Second, RC-LoRA is higher than DD-LoRA in overall ACL Bench accuracy for all three checkpoints: 31.7\% to 38.0\% for DeepSeek, 54.0\% to 58.7\% for Qwen2.5, and 42.3\% to 46.7\% for Llama3.1. The category-level changes from DD-LoRA to RC-LoRA are also non-negative in all nine model--category comparisons, although their magnitudes differ: DeepSeek gains primarily on Counter and Hierarchy Logic, Qwen2.5 improves across all three categories, and Llama3.1 gains on Hierarchy Logic and Pattern Recognition while remaining unchanged on Counter Logic.

This pattern argues against interpreting fine-tuning as a task-independent increase in reasoning ability. A more plausible interpretation is that ordinary structured-task training strengthens some reusable operations but can also induce capability trade-offs, whereas the constraint-enriched regime is associated with a more consistently positive transfer profile across the evaluated checkpoints. We therefore describe the observed effect as \emph{selective structural transfer with constraint-associated cross-checkpoint consistency}. The relevant unit of transfer is not necessarily the source task as a whole, but operations such as ordered comparison, pattern processing, and conditional action selection.

\subsection{Constraint-Induced Relative Preference Shift and the Generation Gap}
\label{sec:discussion_preference}

The constraint-sensitive probes clarify what the current BCDM result does and does not establish. BCDM measures how an explicit prohibition changes the relative log-probability margin between the forbidden action and the strongest legal alternative. The substantially larger mean BCDM for RC-LoRA therefore indicates a larger constraint-induced relative preference shift than in Base, but it should not be interpreted as direct evidence that probability mass is successfully redirected enough to change the model's final decision. The direct-generation diagnostic makes this distinction explicit. Base produces a legal action in 20 of 60 probes, whereas RC-LoRA produces the forbidden action in all 60. The present evidence therefore does not support a claim that RC-LoRA improves executable constraint following under greedy decoding. Instead, it shows a probability-level change in the legal-versus-forbidden margin that remains insufficient to move the generated output across the decision boundary. This preference--generation gap is itself informative: a model can become more responsive to a constraint in relative candidate scores while becoming no more, or even less, compliant at the final-action level.

The constraint-injected examples differ from ordinary state--action supervision because the appropriate supervised target can depend on an additional rule: the same underlying state may require a different target once a preferred option becomes invalid. The paired probes therefore provide a direct behavioral view of how the RC-LoRA model responds when the admissible action space is explicitly changed.

\subsection{Sparse and Non-Additive Causal Reorganization}
\label{sec:discussion_reconfiguration}

The complete Qwen2.5-7B head scan reveals changes in causal dependence unevenly across attention heads. Most heads exhibit relatively small Base-to-RC-LoRA changes, whereas a much smaller subset shows large positive $\Delta CE$ values and some heads show negative changes. This distribution is inconsistent with uniform strengthening of the attention system and instead supports a selective reorganization of functional dependence. L0-H22 provides a representative example. The head contributes positively to the measured behavior in the Base model and becomes substantially more important after fine-tuning. This is consistent with prior findings that fine-tuning can amplify mechanisms already present in a pretrained model rather than necessarily constructing entirely new computational structures \cite{prakash2024finetuning}. The present evidence supports this interpretation for some heads, but not as a universal account of every causal change.

The Top-$k$ experiments provide a stronger test than single-head ranking alone. Joint ablation of the highest-ranked heads produces greater degradation than both global-random and layer-bin-matched random controls, showing that the selected head set carries information beyond intervention size and broad depth-region placement. At the same time, degradation is non-monotonic as $k$ increases. Individual head effects therefore cannot be treated as additive; redundancy, compensation, inhibitory interactions, or larger network-state changes may contribute to the joint-ablation pattern.

The exploratory attention-based screening also differs from the causal results. More generally, prior work cautions against treating attention-based signals as faithful explanations of model behavior \cite{jain2019attention}. The absence of RC-LoRA heads surviving the correlation-based screening therefore does not imply an absence of functionally important heads, because causal ablation identifies several with large behavioral effects. Attention correlation and causal contribution should be treated as distinct quantities.

Several high-$\Delta CE$ heads occur in early layers, but the strongest effects are not confined there. High-impact heads also appear at intermediate depths, and the layer-bin-matched controls show that the Top-$k$ effect is not reproduced by matching only the broad depth profile. However, because exact layers are not matched, generic sensitivity differences among individual layers cannot be fully excluded. The causal redistribution is therefore better characterized as sparse and distributed across depth, with this remaining control limitation made explicit.

\subsection{Scope, Limitations, and Causal Boundary}
\label{sec:limitations}

The behavioral and mechanistic results should be interpreted at different levels. ACL Bench provides evidence of selective cross-task transfer. The constraint probes measure a larger constraint-induced relative preference shift for RC-LoRA than for Base, while the direct-generation diagnostic shows that this does not amount to improved executable compliance. The attention-head interventions establish causal support for the BCDM-based preference shift within the Base-to-RC comparison, but the current experiments do not directly show that the same heads are responsible for ACL Bench gains. This distinction is important because the mechanistic analysis is evaluated with BCDM rather than ACL Bench. A direct causal bridge would require evaluating ACL Bench under the same Top-$k$ and appropriately matched interventions. Stronger ACL degradation under Top-$k$ ablation would support a shared internal basis between constraint adaptation and transfer; stable ACL performance would instead suggest partially distinct mechanisms.

Several additional limitations remain. Dou Dizhu is the only source domain, and ACL Bench is intentionally constructed around structurally related operations, so broader generalization requires additional source domains and independent benchmarks. Complete causal analysis is currently available only for Qwen2.5-7B, and the present attention-head interventions characterize a selected set of functionally important components rather than a complete circuit. The Top-$k$ experiment is also an in-sample validation because the same 60 probes are used both to rank heads by $\Delta CE$ and to evaluate joint ablation. In addition, each random-control mean is based on five sampled head sets, so the controls should be interpreted descriptively rather than as a full randomization distribution. Top-$k$ ablation therefore identifies a functionally important subset for the present probe set, not a complete circuit or a held-out mechanistic signature.

The most informative extensions are held-out or cross-validated Top-$k$ validation, direct ACL Bench ablation with Top-$k$ and corresponding controls, causal replication on additional checkpoints, and evaluation on additional source domains and independent benchmarks. Interaction analysis among the highest-ranked heads would also help distinguish redundancy, compensation, and other sources of the observed non-monotonic joint effects.



